%% ma: L10-B14-0094
%% lop: 10
%% chuong: 5
%% bai: 14
%% dang: 10.14.12
%% loai: DS
%% muc: [NB, NB, TH, TH]
%% dap_an: "ĐSĐĐ"
%% nguon: "chuyen-de-cac-so-dac-trung-cua-mau-so-lieu-khong-ghep-nhom-toan-10-chuong-trinh-moi.pdf (GV. Phan Nhật Linh) (P05), Bài 03 Dạng 2 Phần II Câu 11"
%% kien_thuc: "Lớp 10 Bài 13 (số trung bình, trung vị, tứ phân vị, mốt), Bài 14 (khoảng biến thiên, khoảng tứ phân vị, phương sai, độ lệch chuẩn, giá trị bất thường)"
%% trang_thai: chinh-thuc
%% ghi_chu: "Tự duyệt theo yêu cầu giáo viên 22/10/2026."
\begin{ex}
Người ta thống kê số gia cầm bị tiêu hủy trong vùng dịch của $6$ xã A, B, C, D, E, F như sau (đơn vị: nghìn con):
\begin{center}
\begin{tabular}{|c|c|c|c|c|c|c|}\hline
Xã & A & B & C & D & E & F\\\hline
Số gia cầm bị tiêu hủy & $12$ & $25$ & $27$ & $15$ & $45$ & $5$\\\hline
\end{tabular}
\end{center}
Xét tính đúng sai của các khẳng định sau:
\choiceTF{\True Số trung bình của mẫu số liệu là $21{,}5$}{Khoảng biến thiên của mẫu số liệu là $25$}{\True Trung vị của mẫu số liệu là $20$}{\True Độ lệch chuẩn của mẫu số liệu (làm tròn đến hàng phần trăm) là $s\approx12{,}91$}
\loigiai{Sắp xếp: $5$; $12$; $15$; $25$; $27$; $45$.

a) Đúng: $\overline x=\dfrac{129}6=21{,}5$.

b) Sai: $R=45-5=40$.

c) Đúng: $M_e=\dfrac{15+25}2=20$.

d) Đúng: $s^2\approx166{,}58$; $s\approx12{,}91$.}
\end{ex}
